\section{Results and Discussion}
\label{sec:results_discussion}

\subsection{Physical Structure of the Shadow Boundary}
The numerical and analytical results established in this work provide a comprehensive, unified description of black hole shadows across the entire Kerr parameter space:
\begin{enumerate}
    \item \textbf{Schwarzschild Baseline:} In the absence of rotation ($a = 0$), spherical symmetry enforces an exact circular shadow silhouette of radius $R = 3\sqrt{3}M \approx 5.196M$ on the observer's screen, independent of observer inclination $\theta_o$. Every point along this circular boundary maps to the unstable spherical photon orbit at the photon sphere $r = 3M$.
    \item \textbf{Frame Dragging and Asymmetry:} As black hole spin increases ($0 < |a| < M$), frame dragging breaks spherical symmetry. For an equatorial observer ($\theta_o = 90^\circ$), co-rotating (prograde) photons experience a reduced centrifugal barrier and can approach closer to the black hole before escaping ($r_{\rm ph}^+ < 3M$), resulting in a smaller impact parameter magnitude. Conversely, counter-rotating (retrograde) photons encounter an enhanced centrifugal barrier, requiring a larger impact parameter to orbit ($r_{\rm ph}^- > 3M$). Under Bardeen's celestial screen convention $\alpha = -\xi/\sin\theta_o$, this asymmetry displaces the shadow centroid toward the retrograde side ($\alpha > 0, \bar{\alpha} > 0$) while compressing the prograde boundary.
    \item \textbf{Invariance of Equatorial Vertical Diameter and Horizontal Compression:} At equatorial inclination ($\theta_o = 90^\circ$), the vertical diameter is rigorously invariant at $D_{\mathrm{V}} = 6\sqrt{3}M \approx 10.392M$ across all spins $a \in [0, M]$. For rotating black holes ($0 < a \le M$), $\beta(r) = \pm\sqrt{\eta(r)}$ attains its global stationary maximum where $\eta'(3M) = 0$ with $\eta(3M) = 27M^2$ (at horizontal coordinate $\alpha = +2a$). For Schwarzschild ($a = 0$), $D_{\mathrm{V}} = 6\sqrt{3}M$ follows independently from spherical circular symmetry, with $\alpha = +2a \to 0$ smoothly restoring the unshifted profile. Consequently, rotational deformation at $\theta_o = 90^\circ$ is manifest purely as asymmetric horizontal compression of the prograde edge ($D_{\mathrm{H}} < D_{\mathrm{V}}$). At non-equatorial inclinations, the silhouette smoothly contracts toward the polar circular radius $R_{\rm pole}(a) = \sqrt{\eta(r_p) + a^2} < 3\sqrt{3}M$.
    \item \textbf{Equatorial D-Shape and NHEK Horizon Exposure:} For extreme spin ($a/M \to 1.0$) observed at supercritical inclinations ($\theta_o > \theta_{\rm crit} \approx 47.0586^\circ$), the prograde boundary flattens into a vertical line segment at $\alpha = -2M/\sin\theta_o$. As proven by Bardeen \cite{Bardeen1973} and clarified in modern holographic contexts \cite{Gralla2020}, this flat segment represents the projection of the infinite-throat Near-Horizon Extreme Kerr (NHEK) geometry, where a continuous family of bound spherical photon orbits resides at the degenerate horizon $r = M$.
\end{enumerate}

\subsection{Near-Polar Symmetry and First-Order Invariance}
Our near-polar perturbation analysis reveals an elegant geometric cancellation: although both screen coordinates $\alpha(\phi)$ and $\beta(\phi)$ undergo $\mathcal{O}(\sin\theta_o)$ deformations ($s\,\delta_0 - s\,\lambda\cos(2\phi)$ and $s\,\lambda\sin(2\phi)$, respectively), the metric speed $\sigma(\phi) = R_p + 2s\lambda\cos\phi + \mathcal{O}(s^2)$ is a pure $\cos\phi$ dipole, while the planar curvature magnitude $\kappa(\phi) = R_p^{-1} + \mathcal{O}(s^2)$ is invariant to first order.

Consequently, by dipole orthogonality ($\int_0^{2\pi}\cos\phi\,d\phi = 0$), both the continuous boundary perimeter $P_{\rm cont}$ and the curvature functional $\mathcal{I} = \int_0^{2\pi} \kappa^2 \sigma^3\,d\phi$ remain invariant to first order in inclination:
\begin{equation}
    P_{\rm cont}(\theta_o) = 2\pi R_p + \mathcal{O}(\sin^2\theta_o), \qquad \mathcal{I}(\theta_o) = 2\pi R_p + \mathcal{O}(\sin^2\theta_o).
\end{equation}
Physically, this indicates that for observers viewing a rotating black hole within $\approx 5^\circ$ of its rotation axis, the shadow perimeter and optical cross-section are remarkably insensitive to spin-induced deformations, deviating from a perfect circle only at second order in inclination. For larger viewing angles, such as the astrophysical inclination inferred for the M87* relativistic jet ($\theta_o \approx 17^\circ$, where $s \equiv \sin\theta_o \approx 0.29$), the first-order asymptotic regime serves as a valuable analytical baseline, while quadratic corrections $\mathcal{O}(s^2)$ enter at the $\sim 5\text{--}8\%$ level and must be accounted for via the full non-linear parameterization.

\subsection{Relevance to Current and Future Astrophysical Observations}
Direct observations of M87* \cite{EHT2019_PaperI} and Sgr A* \cite{EHT2022_PaperI} by the Event Horizon Telescope have measured the angular diameter of the central brightness depression to within $\approx 10\%$ uncertainty, demonstrating consistency with the predictions of general relativity for Kerr black holes.

Future high-frequency space-based very-long-baseline interferometry (VLBI) missions (such as the next-generation EHT [ngEHT] and orbital space-VLBI concepts) aim to resolve the sharp, exponentially demagnified sub-rings of the photon ring \cite{Johnson2020,Chael2021,Gralla2020}. Because the critical boundary curves constructed by our framework define the mathematical envelope to which these photon sub-rings asymptotically converge, high-precision, scale-invariant geodesic integration frameworks such as the one presented here will be essential for extracting black hole mass, spin, and inclination from future precision interferometric closure phases.

\subsection{Methodological Distinction: Analytical Shadow Construction vs.\ Numerical Integration}
It is important to emphasize a key methodological distinction within this work:
\begin{itemize}
    \item \textbf{Analytical Shadow Construction:} The black-hole shadow boundaries in Section~\ref{sec:shadow_construction} and their observables in Table~\ref{tab:shadow_observables} are constructed directly from the analytical parameterization of spherical photon orbits $(\xi(r), \eta(r))$ via the offset variable $u = r - 3M$ and Brent root-finding. This approach guarantees machine-precision boundaries without ray-tracing noise.
    \item \textbf{Numerical Geodesic Integration:} The DOP853 numerical ODE integrator in Section~\ref{sec:numerical_architecture} is a general-purpose, continuous 8th-order integration engine developed to propagate arbitrary timelike and null trajectories in Kerr spacetime, verified against conserved quantities ($|\Delta\mathcal{H}_{\rm norm}| < 10^{-10}$), physical epicyclic frequencies, and logarithmic dwelling scaling near separatrices.
\end{itemize}
These two complementary components provide independent mutual verification: the numerical integrator validates trajectory dynamics in the strong-field regime, while the analytical shadow curves establish exact ground-truth geometric benchmarks.
